\documentclass[11pt]{article}

\usepackage[margin=1in]{geometry}
\usepackage{booktabs}
\usepackage{hyperref}
\usepackage{enumitem}
\usepackage{array}
\usepackage{xurl}

\title{Thin Harness, Strong Contracts:\\Production-Oriented Agent Harnesses for Stateful AI Agents}
\author{Song Luo\\\texttt{luosongred@gmail.com}}
\date{June 2026}

\begin{document}
\maketitle

\begin{abstract}
Language-model agents are increasingly expected to operate software, browse information sources, update files, schedule jobs, manage memory, and prepare publishable artifacts. In such settings, final-answer evaluation is insufficient: an agent may produce a plausible answer while reading inappropriate context, mutating the wrong state, bypassing review, or making an irreversible tool call. This paper argues for a production-oriented view of agent harnesses: not merely benchmark runners, but controlled execution and evaluation environments that record trajectories, enforce permissions, compare state diffs, support replay, and preserve auditability. Drawing on a first-party artifact base of skill, memory, security, scheduler, and AI-radar systems, we distinguish \emph{capability-patch harnesses}, which compensate for temporary model weaknesses and often become technical debt, from \emph{system-boundary harnesses}, which encode durable trust boundaries. We propose the design principle \emph{thin harness, strong contracts}: keep orchestration thin enough to avoid fighting model progress, but make contracts for tools, memory, side effects, review, and replay explicit. A 24-case pilot benchmark across trigger routing, stateful tool use, permission security, memory scoping, and replay illustrates how such contracts can be evaluated. The contribution is a bounded engineering framework for making stateful agents more controllable and inspectable under production constraints.
\end{abstract}

\section{Introduction}

The practical unit of work for an agent is no longer always a text answer. A coding agent edits a repository and runs tests. A research agent reads sources, builds a report candidate, and preserves citations. A scheduling agent creates or edits future jobs. A memory agent decides what should be recalled, updated, or forgotten. In each case, the important question is not only whether the final response sounds correct, but whether the agent changed the world in the right way.

This creates a gap between conventional language-model evaluation and production agent operation. Outcome metrics can hide unsafe trajectories. An agent may answer correctly after reading a secret file; it may generate a useful summary after merging stale memory from another project; it may complete a task by sending an email that should have required human approval. In production, these are failures even when the final answer is acceptable.

This paper uses \emph{agent harness} to mean the controlled environment around an agent: the runtime, tools, state, policies, logs, replay mechanism, evaluation checks, and review loop that make agent behavior observable and governable. This is broader than a benchmark runner and narrower than a full product. The harness is the layer that answers questions such as: What did the agent do? Which state changed? Was each tool call allowed? Can the failure be replayed? Who approved the side effect? What changed between two model or prompt versions?

The central claim is that harnesses should be judged by whether they patch a temporary model weakness or define a durable system boundary. Thick wrappers that hard-code planning, repair, routing, or role choreography around one generation of models can become technical debt. By contrast, permission boundaries, state diffs, audit logs, replayable traces, safety cases, and human review remain necessary as models improve. Stronger models can do more; therefore the need for trustworthy boundaries increases rather than disappears.

\section{Background}

\subsection{From Prompts to Operational Skills}

Reusable skills provide a compact setting for studying harness design. In practice, a skill is not only an instruction prompt. It may include trigger metadata, input requirements, tool permissions, execution steps, output contracts, memory rules, tests, and failure handling. Treating a skill as prose encourages over-broad activation and unstable outputs. Treating it as an interface makes the relevant questions explicit: when should it run, what context may it consume, what may it change, what should it return, and when must it stop.

The first-party SkillOps artifact base frames this as an operational lifecycle: specify, lint, test, operate, audit, and retire skills. The same pattern extends to broader agents. A production harness makes the agent's interface with the world explicit.

\subsection{Related Systems}

Recent work has moved agent evaluation beyond static input-output tests. ToolSandbox emphasizes stateful, conversational, interactive evaluation with intermediate milestones over arbitrary trajectories~\cite{toolsandbox}. AgentDojo evaluates prompt-injection attacks and defenses in dynamic tool-calling environments containing untrusted data~\cite{agentdojo}. SWE-agent argues that agent-computer interfaces affect agent behavior and performance in software engineering tasks~\cite{sweagent,aci}. Deli AutoResearch contributes a long-horizon operating protocol based on persistent state files, stall detection, fresh sessions, heartbeat monitoring, and independent verification~\cite{deli}.

These systems motivate the present focus but do not fully replace it. This paper asks how production teams should design the harness boundary itself: the contracts for state, tools, permissions, memory, replay, audit, and review that sit between a capable model and a real workflow.

\subsection{Evidence Sources}

The empirical basis for this paper has four layers. First, local essays and notes provide the conceptual vocabulary around routing contracts, skill interfaces, memory boundaries, and thin-versus-thick harnesses. Second, a generated GitHub snapshot records metadata for first-party skill, memory, security, self-audit, and AI-radar repositories, plus upstream OpenClaw and Hermes pull requests used as boundary examples. Third, a source evidence audit records path-level and keyword-signal evidence for first-party repositories without storing source bodies. Fourth, the pilot benchmark contributes raw JSONL runs, regenerated tables, failure inspection, and prompt-hash metadata for live model calls. This evidence structure is intentionally conservative: repository metadata supports artifact existence and scope; source-signal snapshots support bounded artifact interpretation; only raw run files support reported pilot metrics.

\section{A Taxonomy of Harnesses}

\subsection{Capability-Patch Harnesses}

Capability-patch harnesses compensate for weaknesses of a current model. Examples include rigid planners, brittle JSON repair loops, role-heavy multi-agent scripts, hard-coded tool routers, and checklists that repeat steps the model should eventually learn to perform. These layers can be useful, but they are coupled to a particular model generation. When the model's native planning, tool use, long-context handling, or structured output improves, the patch can start to constrain the system.

The risk is not that all orchestration is bad. The risk is that a thick harness silently becomes the product's real logic while remaining harder to test than ordinary code and less flexible than the model it wraps.

\subsection{System-Boundary Harnesses}

System-boundary harnesses define the interface between agent cognition and external responsibility. They do not primarily tell the model how to think; they define what it may access, mutate, publish, remember, forget, and escalate. Examples include:

\begin{itemize}[leftmargin=*]
  \item tool-call logging and argument capture;
  \item state snapshots and final state diffs;
  \item dry-run previews before side effects;
  \item permission gates for file, database, email, browser, and scheduler actions;
  \item memory scoping and forgetting policies;
  \item prompt-injection and secret-leakage regression cases;
  \item replayable traces for debugging;
  \item human review queues and audit logs.
\end{itemize}

These needs do not disappear with better models. They are part of the trust and accountability structure around any system that acts.

\section{Thin Harness, Strong Contracts}

The proposed design principle is \emph{thin harness, strong contracts}. The harness should stay thin where it would otherwise duplicate model cognition, but strong where it defines external effects. In practical terms, this means:

\begin{enumerate}[leftmargin=*]
  \item Evaluate trajectories, not only final answers.
  \item Treat state diffs as evidence of completion.
  \item Make tool permissions explicit and testable.
  \item Separate stateless capabilities from stateful memory.
  \item Use dry-run and review states before irreversible actions.
  \item Keep replay artifacts sufficient to localize failures.
  \item Version prompts, tool schemas, and harness policies together.
\end{enumerate}

This principle also reframes the skill-memory distinction. Skills are reusable procedural capabilities; memory is bounded state. The former can be versioned and shared. The latter must remain scoped, erasable, and protected against cross-project contamination. In a richer agent system, repository rules, skills, memories, tools, subagents, and external connectors occupy different layers of the harness.

\section{Case Studies}

\subsection{SkillOps}

SkillOps is used here as an artifact-based example of treating skills as operational artifacts rather than prompt snippets. The inventory, repository metadata, and source-signal audit associate it with benchmark inputs, experiment runners, results, and tests for modular skills. The design lesson is therefore bounded: skill quality should be assessed by routing, input contracts, output structure, failure paths, and tests under noisy inputs and boundary cases, rather than by one impressive demo.

\subsection{AI Radar}

The AI Radar web and skill artifacts, as recorded in the artifact inventory and source-signal audit, illustrate workflow-level harnessing as a design pattern. The repository metadata, path-level signals, and local notes describe source ingestion, understanding, retrieval, report candidate generation, review, audit, and publication boundaries. This is not a benchmark harness, but it motivates the same production pattern: each step should expose state, evidence, permissions, and review semantics.

\subsection{OpenClaw and Hermes PRs}

Several upstream PRs recorded in the inventory provide concrete boundary examples at the metadata and change-description level. OpenClaw cron dry-run work is treated as evidence for preview-before-mutation semantics. Timezone and time-only scheduling PRs are treated as evidence that scheduler semantics belong in the harness contract. Hermes reusable-session and scoped-recall PRs are treated as evidence for continuity and memory-isolation concerns in recurring agent work.

\subsection{Security, Memory, and Self-Audit Skills}

The security, memory, and self-audit repositories are used as artifact examples for three recurring harness concerns: preflight scanning, bounded memory, and periodic health checks. The inventory and source-signal audit associate skill-security-guard with tests and risk dimensions, persistent-memory with memory-scoping files, and agent-self-audit with audit criteria. These examples support the taxonomy but should not be read as runtime proof.

\section{Pilot Benchmark}

We constructed 24 pilot cases across five categories: trigger routing, stateful tool use, permission/security, memory scoping, and replay/recovery. Each case specifies an initial state, allowed tools, expected required tools, forbidden tools, and state assertions. The harness evaluates action plans by simulating tool effects, counting invalid tool calls and permission violations, and checking the final state.

Three prompt variants are supported. The \texttt{no\_harness} variant gives the task with minimal structure. The \texttt{thick\_checklist} variant adds broad safety and quality reminders. The \texttt{thin\_contract} variant gives a compact tool contract, forbidden side effects, untrusted-data rule, and strict JSON schema. This setup is intentionally small; its purpose is to demonstrate an evaluation shape rather than claim general model performance.

\begin{table}[h]
\centering
\begin{tabular}{llrrrrrr}
\toprule
Provider & Variant & Total & Executed & Provider errors & Parse errors & Success & Violations \\
\midrule
codex\_reference & no harness & 24 & 24 & 0 & 0 & 1.000 & 0 \\
codex\_reference & thick checklist & 24 & 24 & 0 & 0 & 1.000 & 0 \\
codex\_reference & thin contract & 24 & 24 & 0 & 0 & 1.000 & 0 \\
deepseek & no harness & 24 & 24 & 0 & 0 & 0.000 & 0 \\
deepseek & thick checklist & 24 & 24 & 0 & 0 & 0.708 & 0 \\
deepseek & thin contract & 24 & 24 & 0 & 0 & 0.583 & 0 \\
kimi & thin contract & 1 & 0 & 1 & 0 & 0.000 & 0 \\
\bottomrule
\end{tabular}

\caption{Pilot results regenerated from raw JSONL. The reference policy checks benchmark mechanics. DeepSeek rows are live API outputs. The Kimi row records an authentication failure and is not a model-behavior measurement.}
\label{tab:pilot}
\end{table}

The deterministic reference policy is not a model benchmark. It verifies that the case definitions, simulator, metrics, and aggregation logic are internally consistent. The live DeepSeek run contains 72 model-backed rows. Three rows carry \texttt{usage.rerun\_reason=parse\_error} after parser-instrumentation repair; two rows remain classified as parse errors. Parser failures are classified separately from provider availability failures.

In this pilot, the unstructured no-harness prompt failed all cases under the action-plan scoring because it usually did not emit executable tool actions. The thick checklist variant achieved a higher success rate than the text-only thin-contract prompt. Appendix Tables~\ref{tab:category-summary} and~\ref{tab:failure-taxonomy} show that the remaining failures are dominated by routing and state-diff errors, not permission violations. This does not refute the strong-contract argument; rather, it sharpens it. A contract written only as natural language is still a prompt. A production harness needs enforceable schemas, tool APIs, parsers, retries, and state checks around that contract.

The Kimi provider attempt returned authentication failure at the official OpenAI-compatible endpoint, so no Kimi model-behavior result is reported. All live provider runners read credentials only from process environment variables~\cite{deepseekapi,kimiapi}.

\section{Discussion}

The taxonomy clarifies the common disagreement about whether harnesses become technical debt. Both positions are partly correct. Harnesses that encode temporary cognitive scaffolding can become debt. Harnesses that encode trust boundaries become infrastructure. Model progress may absorb planning, tool selection, and format-following behaviors, but it does not remove the need to know what state changed, what permissions were exercised, which memory was used, or who approved an external side effect.

The pilot results also warn against a superficial reading of ``thin.'' Thin does not mean weak, informal, or merely short. A thin contract that is only described in natural language can fail at routing or formatting. The durable part of the harness is the enforceable boundary: schemas, parsers, permission checks, state assertions, replay logs, and review gates. This is why the practical target is not fewer controls, but less brittle cognitive choreography and stronger external contracts.

This also explains why model laboratories often produce important agent-harness work. Teams close to model training and failure analysis see the hooks earlier: where the model loses state, confuses instructions with data, overuses tools, ignores uncertainty, or fails to recover. Harnesses then become part of the feedback loop: trajectory traces, state diffs, safety failures, and review labels can inform evaluation, post-training, and product design.

\section{Limitations}

The artifact base is first-party and author-interpreted. Several case-study claims are supported by generated repository and pull-request metadata rather than by vendored source snapshots in this repository. Repository metadata supports claims about artifact existence, ownership status, broad scope, and PR shape; it does not by itself prove all implementation behavior.

The benchmark is small, synthetic, and designed for clarity rather than statistical generality. The deterministic reference runner checks harness mechanics, not model capability. Live API results are subject to model-version, endpoint, authentication, and parameter changes. The case studies are engineering evidence, not controlled deployments across independent organizations. These limitations make the claims bounded: the paper proposes a design framework and pilot evaluation shape, not a universal benchmark.

\section{Conclusion}

Production agents need more than better final answers. They need controlled interfaces to state, tools, memory, and people. The useful harness is not a thick substitute for model intelligence. It is a thin but firm boundary around action: explicit contracts, replayable trajectories, state diffs, permissions, audit, and review. As models improve, brittle capability patches should shrink. Strong system-boundary harnesses should become the infrastructure that lets stronger agents be trusted in real workflows.

\appendix

\section{Benchmark Case Design}

The pilot benchmark is intentionally compact. Its role is to make the harness contract inspectable, not to provide a general-purpose agent leaderboard. Table~\ref{tab:cases} summarizes the case categories.

\begin{table}[h]
\centering
\begin{tabular}{lrl}
\toprule
Category & Cases & Boundary tested \\
\midrule
Trigger routing & 8 & Whether the skill should activate, refuse, or request review \\
Stateful tool use & 6 & Whether files, tests, and report states change as expected \\
Permission/security & 4 & Whether unsafe tool use, email, publication, and secrets are blocked \\
Memory scoping & 3 & Whether memory is searched, forgotten, or refused within scope \\
Replay/recovery & 3 & Whether recovery actions and replayable traces are preserved \\
\bottomrule
\end{tabular}
\caption{Pilot case distribution.}
\label{tab:cases}
\end{table}

Each case records a task, an initial state, allowed tools, expected required tools, forbidden tools, optional untrusted observations, and state assertions. The simulator then evaluates whether the proposed action plan changes the mock world in the expected way.

\section{Artifact Checklist}

The repository contains the artifacts needed to audit the pilot:

\begin{itemize}[leftmargin=*]
  \item benchmark cases in \texttt{benchmark/cases.jsonl};
  \item harness and metric logic in \texttt{harness/};
  \item live and deterministic raw rows in \texttt{experiments/raw/};
  \item regenerated tables in \texttt{results/tables/};
  \item generated GitHub metadata snapshots in \texttt{evidence/};
  \item path-level source-signal audit in \texttt{evidence/source\_audit.md};
  \item reproduction and claim-boundary notes in \texttt{docs/}.
\end{itemize}

The raw rows are the source of reported pilot metrics. GitHub metadata snapshots support artifact and PR identification. The source audit supports bounded path-level interpretation, not implementation-behavior claims beyond the case-study framing.

\section{Additional Pilot Tables}

\begin{table}[h]
\centering
\begin{tabular}{llrrrr}
\toprule
Variant & Category & Total & Executed & Success & Routing errors \\
\midrule
no harness & memory scoping & 3 & 3 & 0.000 & 3 \\
no harness & permission security & 4 & 4 & 0.000 & 4 \\
no harness & replay recovery & 3 & 3 & 0.000 & 3 \\
no harness & stateful tool use & 6 & 6 & 0.000 & 6 \\
no harness & trigger routing & 8 & 8 & 0.000 & 8 \\
thick checklist & memory scoping & 3 & 3 & 1.000 & 0 \\
thick checklist & permission security & 4 & 4 & 0.500 & 2 \\
thick checklist & replay recovery & 3 & 3 & 0.000 & 3 \\
thick checklist & stateful tool use & 6 & 6 & 0.333 & 3 \\
thick checklist & trigger routing & 8 & 8 & 0.500 & 4 \\
thin contract & memory scoping & 3 & 3 & 1.000 & 0 \\
thin contract & permission security & 4 & 4 & 0.750 & 1 \\
thin contract & replay recovery & 3 & 3 & 0.000 & 2 \\
thin contract & stateful tool use & 6 & 6 & 0.333 & 2 \\
thin contract & trigger routing & 8 & 8 & 0.625 & 3 \\
\bottomrule
\end{tabular}

\caption{DeepSeek pilot results by prompt variant and case category.}
\label{tab:category-summary}
\end{table}

\begin{table}[h]
\centering
\begin{tabular}{lllr}
\toprule
Provider & Variant & Failure class & Count \\
\midrule
deepseek & no harness & parse error & 2 \\
deepseek & no harness & routing error & 24 \\
deepseek & no harness & state diff error & 10 \\
deepseek & thick checklist & routing error & 7 \\
deepseek & thick checklist & state diff error & 6 \\
deepseek & thin contract & routing error & 9 \\
deepseek & thin contract & state diff error & 5 \\
kimi & thin contract & provider error & 1 \\
kimi & thin contract & routing error & 1 \\
\bottomrule
\end{tabular}

\caption{Failure taxonomy regenerated from raw rows. A single row may contribute to more than one failure class.}
\label{tab:failure-taxonomy}
\end{table}

\bibliographystyle{plain}
\bibliography{references}

\end{document}
